\section{Finite-state filtering after every report}\label{sec:streaming}
Fix $N\ge2$ and abbreviate $d_n=d_A(n,N-n)$ and $D_*=D_A(N)>0$.
The exact state from Theorem~\ref{thm:causal} will be denoted $s_n$.
A clock $n$, calibration, prior moments and a precomputed finite program
are read-only data. The charged resource is the persistent state between
successive input symbols.

The resource model and loss are defined in Section~\ref{sec:experiment}.

\begin{lemma}[Uniform finite-horizon transition bounds]\label{lem:lipschitz}
On the compact exact state sets there are constants $L_n<\infty$ such
that the update $T_n(s,g,x)$ satisfies
\[
 \norm{T_n(s,g,x)-T_n(\widetilde s,g,x)}
 \le L_n\norm{s-\widetilde s}
\]
for every admitted $g,x$ and $s,\widetilde s\in\cS_n$.
Every fixed finite probe prediction vector is likewise Lipschitz in the
state, with a finite constant independent of the history.
\end{lemma}
\begin{proof}
In factor coordinates an update before changeover appends the new factor
$f/\mu f$. The old factor coordinates are unchanged. All unnormalized
report probabilities lie in $[\kappa,1]$; consequently the normalized
factors lie pointwise in $[\kappa,1/\kappa]$. Their coordinates are bounded,
since they are continuous images of a compact cube and finite alphabet.
At changeover each needed moment is a ratio of polynomial expressions in
these coordinates and the fixed prior moments, with denominator the
integral of a product of normalized factors, bounded below by $\kappa^{n+1}$.
The same lower bound holds on line segments between factor tuples,
because each individual interpolated normalized factor is still at least
$\kappa$. Finite differentiation and the mean value theorem give a
uniform derivative bound there.

In moment coordinates the update is \eqref{eq:sparse-update}. Its
denominator is at least $\kappa$ at every state; on a line segment between
two such moment vectors the corresponding posterior mixture has the same
lower bound. Numerators, coefficients and moment coordinates range over
compact sets. The quotient rule therefore gives a uniform bound. Probe
predictions are linear in the relevant moment coordinates, or normalized
polynomial expressions in the factor coordinates with the same positive
denominator control. The argument gives their claimed Lipschitz constants.
All constants use only finite algebra and the supplied moments; they may
grow with $N$ or deteriorate with $\kappa$.
\end{proof}

\begin{theorem}[Finite-state streaming upper bound]\label{thm:stream-upper}
For every $M\ge1$ there is a deterministic $M$-state filter, updated after
every report, whose prediction regret at every checkpoint and every
admitted command-report history is at most
\begin{equation}\label{eq:stream-upper}
 C_N M^{-2/D_*},
\end{equation}
where $C_N<\infty$ depends on the fixed experiment, prior, probe lists
and total horizon, but not on $M$ or the input history. The representatives
and read-only transition functions are supplied existentially from the exact
model; this is not an effective synthesis theorem from finite-precision data.
\end{theorem}
\begin{proof}
Enclose $\cS_n\subset\R^{d_n}$ in $[-B_n,B_n]^{d_n}$. For $d_n>0$
partition each coordinate interval into
$k_n=\lfloor M^{1/d_n}\rfloor$ intervals. In each nonempty cell choose
one representative \emph{in} $\cS_n$. There are at most $k_n^{d_n}\le M$
representatives, and assigning a state to its cell representative has
error at most
\[
 \delta_n\le 2B_n\sqrt{d_n}/k_n
 \le4B_n\sqrt{d_n}\,M^{-1/d_n}
 \le C M^{-1/D_*}.
\]
Here $C$ is the maximum of the finite prefactors. Empty cells need no
representative. For $d_n=0$ use the unique state. Deterministic boundary
and tie conventions give Borel quantizers $Q_n$.

The filter stores only the index of a representative $\widehat s_n$.
Its transition is exactly
\begin{equation}\label{eq:transducer}
 \widehat s_{n+1}
   =Q_{n+1}\bigl(T_n(\widehat s_n,g_{n+1},x_{n+1})\bigr).
\end{equation}
Since each representative is reachable and every report has positive
probability, its exact image lies in $\cS_{n+1}$. Thus the quantizer is
always applied inside its domain and no infeasible moment update is used.
The finite representative choices and the functions in
\eqref{eq:transducer} are read-only program data. They may be obtained by
selecting points in nonempty cells of the compact state sets; no efficient global synthesis
or computability of unsupplied real prior moments is asserted.

Let $e_n=\norm{s_n-\widehat s_n}$. On the actual command-report stream,
Lemma~\ref{lem:lipschitz} gives
\[
 e_0=0,\qquad e_{n+1}\le L_ne_n+C M^{-1/D_*}.
\]
Set $b_0=0$ and $b_{n+1}=L_nb_n+C$. Induction proves
$e_n\le b_nM^{-1/D_*}$. This explicitly accumulates the errors from
\emph{all} previous lossy updates. At a query use the probe predictions
of $\widehat s_n$, which are actual probabilities. If the vector prediction
map has Lipschitz constant $K_n$, \eqref{eq:brier} bounds regret by
$(K_n^2/s_n^{\rm q})b_n^2 M^{-2/D_*}$. The maximum of these finitely
many constants proves \eqref{eq:stream-upper}. No exact prefix was kept
for the changeover: the factor-to-moment computation uses the representative
factors already encoded by $i_n$.
\end{proof}

\begin{theorem}[Matching exponent in one fixed streaming experiment]\label{thm:stream-lower}
Choose $n_*$ attaining $D_*=d_A(n_*,N-n_*)$. For the first $n_*$ trials
supply independent uniform commands from $\cG$, independent of the
parameter, cartridges and filtering coins. The realized command is read
once by the filter. At $n_*$ reveal the independent probe query and execute
its remaining trials. This is one fixed experiment, known before the
filter is chosen. Its optimal unconditional regret over $M$-state filters
satisfies
\begin{equation}\label{eq:stream-match}
 0<c_NM^{-2/D_*}\le\overline{\mathcal R}_{M,N}^{\rm stream}
                  \le C_NM^{-2/D_*}.
\end{equation}
The same power is sharp for the minimax streaming criterion taking the
maximum conditional regret over checkpoints and input histories. Fresh
private update and decoder randomization does not defeat the lower bound;
a fresh public coding seed independent of the complete prefix may also be
allowed. Thus for fixed $N$, the required persistent bits for this prediction
task satisfy
$B_*(\varepsilon)=\tfrac{D_*}{2}\log_2(1/\varepsilon)+O_N(1)$.
\end{theorem}
\begin{proof}
Write $n=n_*$, $d=D_*$, and $s=s_n^{\rm q}$. Theorem~\ref{thm:rank}
gives an interior command tuple $g_0$ where the probe map $p(g)$ on
all-failure prefixes has rank $d$. Select $d$ independent probe coordinates,
and let $E$ denote their coordinate projection, of operator norm one.
Choose complementary command coordinates $w$ so that
$g\mapsto(Ep(g),w(g))$ has an invertible derivative at $g_0$.
Let $\Phi$ be its local inverse. There is a compact product
$B_\rho^d\times W$ inside its domain, centered at the projected
prediction point, such that $\Phi$ stays in the command interior and
$|\det D\Phi|\ge j_0>0$. The complementary box $W$ has positive
volume, with unit-volume convention in zero dimensions.

If $k=Kn$ is the number of real command entries, their density is
$(1-2\eta)^{-k}$. The joint density of commands and an all-failure
prefix is this density times
$Z(g)=\mu\prod_iF_{g_i}\ge\eta^n$. Change variables with $\Phi$ and
integrate over $W$. The subprobability law of the projected predictions
on this prefix dominates $\beta$ times uniform probability on
$B_\rho^d$, with
\[
 \beta=\frac{\eta^n j_0|W|\,|B_\rho^d|}{(1-2\eta)^k}>0.
\]
In particular the lower bound retains the probability of the failure
prefix. It does not assign an atom to an exact real command or assume
a flat ball in the entire curved prediction image.

Every streaming filter is in particular an $M$-message encoder at this
checkpoint. Given any message, average decoder coins conditional on each
query; squared loss can only decrease. This gives at most $M$ prediction
vectors. Randomizing the encoder cannot beat the closest of these vectors
at a fixed history. Projecting their centers by $E$ cannot increase squared
distance. For a uniform point $X$ in a $d$-ball of radius $\rho$ and any
$M$ centers, the union-of-balls estimate gives
\[
 \Pr\{\min_i\norm{X-z_i}\le r\}\le
                    \min\{1,M(r/\rho)^d\}.
\]
Integrating its complementary tail yields
\[
 \E\min_i\norm{X-z_i}^2
 \ge\int_0^{\rho^2 M^{-2/d}}
          (1-Mt^{d/2}/\rho^d)\dd t
 =\frac d{d+2}\rho^2 M^{-2/d}.
\]
Multiply by $\beta/s$ to obtain $c_N>0$. Fresh public randomness can be
fixed and then averaged because it is independent of the prefix. Private
online coins may equivalently be drawn on an independent tape for this
proof; fixing them makes a deterministic transducer and the same bound
holds uniformly. That proof device supplies no tape to the implemented
filter. The prescribed exploration is independent of that tape, a condition
that is essential to this particular argument. An old command-generation
seed is not fresh coding randomness and remains charged.

Theorem~\ref{thm:stream-upper} gives the upper bound in the same experiment.
The minimax conditional criterion is at least its unconditional regret,
while the universal upper filter works at all checkpoints. Finally put
$M=2^B$ and solve the two inequalities. The Euclidean covering mechanism
is classical \cite{GrafLuschgy}; the identified sparse dimension and the
actual repeated-update construction are the experiment-specific inputs.
\end{proof}

\begin{theorem}[A separate finite-state control bound]\label{thm:control-bound}
For the same finite detector and horizon, suppose stage rewards
$r_n(g,x)\in[0,1]$ are continuous in $g$ for each report, and the objective
is their expected sum. There is an $M$-state feedback controller whose
loss relative to the full-history optimum is at most
\begin{equation}\label{eq:control-bound}
 C'_N M^{-1/D_*}.
\end{equation}
No matching lower power for arbitrary control tasks is asserted. An
additional finite payoff automaton must be counted as task state, multiplying
the number of states if it is included.
\end{theorem}
\begin{proof}
Exact dynamic programming on $\cS_n$ is valid by
Proposition~\ref{prop:tests}. With terminal value zero, its recursion is
\[
 V_n(s)=\max_{g\in\cG}\sum_{x\in\cX}p_x(s,g)
       \{r_{n+1}(g,x)+V_{n+1}(T_n(s,g,x))\}.
\]
Here $p_x$ is a linear moment functional or a normalized factor expression.
The positive denominator bounds, compact action set and continuous rewards
make the displayed objective continuous, and the maximum is attained.
By backward induction it is uniformly Lipschitz in $s$, since finite sums,
products of bounded Lipschitz functions, and the updates in
Lemma~\ref{lem:lipschitz} preserve that property. The maximum of functions
with a common Lipschitz constant is again Lipschitz. Denote by $H_n$ a
uniform Lipschitz constant of the action objective, before its maximization.

At each representative choose an exact maximizing action and store that
action in the finite output table. At the true state $s_n$, the loss from
using the action optimal at $\widehat s_n$ is at most $2H_ne_n$.
Update the representative by \eqref{eq:transducer}. The true exact state
is computed only for analysis along the controller's actual observations;
the implemented controller still stores only its index. Bellman telescoping
and conditional expectation give total loss at most
$2\sum_{n<N}H_ne_n\le2\sum_{n<N}H_nb_n M^{-1/D_*}$.
This proves \eqref{eq:control-bound}. It uses no continuity of the optimizing
action as a function of state, and it does not differentiate an optimizer.
Its approximation principle is related to \cite{Subramanian}; it is not
claimed as a new general dynamic-programming principle.
\end{proof}

\subsection{Horizon and calibration remain part of the constants}
If every factor in a probe menu satisfies $0<F\le v<1$, then
$H_{n,j}\le v^{N-n}$ and every posterior prediction satisfies the same
bound. The one-state decoder that always announces zero therefore has
\begin{equation}\label{eq:zero-memory}
 \mathcal R\le v^{2(N-n)}.
\end{equation}
For the shared physical menu $F=c+\delta k_j$ one may take $v=c+\delta$.
Thus a checkpoint implementation may choose, in advance, between the
streaming filter and the one-state decoder and obtain
\[
 \mathcal R\le\min\{C_N M^{-2/D_*},(c+\delta)^{2(N-n)}\}.
\]
At $c=1/2$, $\delta=1/8$, $N-n=20$, the second bound is
$(5/8)^{40}<10^{-8}$. This does not contradict a fixed-$N$ high-resolution
exponent. Neither the new sparse rank theorem nor the streaming construction
turns it into a horizon-uniform useful-memory law. We have not rescaled
the loss or discarded rare histories to avoid this bound.

The upper constants involve state bounds, transition Lipschitz constants
and repeated error propagation. The lower constants involve an attainable
Jacobian, a chart radius, evidence and the prior pairing. They are not
uniform over vanishing detector amplitudes, all full-support priors or
increasing horizons. With a finite fixed command alphabet and finite $N$
there are finitely many histories and an exact finite code eventually
suffices; the matching asymptotic lower bound above expressly uses the
fixed continuous-command exploration protocol.
